% !PRINTED 695
\weberpage{695}
%<<CONT
eingeschlossen, und es ist dem absoluten Werthe nach
\[
R_m < \frac{2\,C}{m^s},
\]
oder, wenn $s > c$ ist,
\[
R_m < \frac{2\,C}{m^c}.
\]

Diese obere Grenze für $R_m$, die von $s$ unabhängig ist, kann aber, wenn $c$ positiv ist, dadurch, dass man $m$ hinlänglich gross annimmt, unter jeden noch so kleinen Werth herabgedrückt werden.

Daraus folgt aber nicht nur die Convergenz, sondern auch die Stetigkeit von $S$. Denn nimmt man $m$ hinlänglich gross, so wird nicht nur $R_m$, sondern es werden auch die Schwankung von $R_m$ bei veränderlichem $s$ unendlich klein, und $S_m$ ist für ein feststehendes $m$ eine stetige Function von $s$. Also ist auch $S$ stetig.
% [sic] The source has a singular subject (die Schwankung) with the plural verb werden.

Es ist dabei noch zu bemerken, dass die Voraussetzung über $\sigma_m$ keineswegs die Convergenz der unendlichen Reihe $\Sigma a_k$ voraussetzt. Es ist nicht einmal erforderlich, dass die $a_k$ reell seien; denn wenn sie imaginär sind, so braucht man nur den Satz auf den reellen und den imaginären Bestandtheil anzuwenden. Endlich ist es auch nicht nothwendig, die $a_1, a_2, a_3, \dots$ als Constanten vorauszusetzen. Alles bleibt gültig, wenn es \gesperrt{stetige Functionen} von $s$ sind.

2. \gesperrt{Bedeutet $\mu_1, \mu_2, \dots, \mu_n, \dots$ eine unendliche Menge positiver Zahlen von der Beschaffenheit, dass zwei endliche Zahlen $\alpha, \beta$ so angegeben werden können, dass für jedes noch so grosse $n$}
\weq{3}{\alpha < \dfrac{n}{\mu_n} < \beta,}
\gesperrt{so ist die unendliche Reihe}
\[
S = \frac{1}{\mu_1^s} + \frac{1}{\mu_2^s} + \frac{1}{\mu_3^s} + \cdots
\]
\gesperrt{convergent, so lange der Exponent $s$ grösser als 1 ist, und das Product $(s-1)\,S$ bleibt bei hinlänglicher Annäherung von $s$ an den Werth 1 gleichfalls zwischen den Grenzen $\alpha$ und $\beta$ oder nähert sich wenigstens einem dieser Werthe unbegrenzt an.}
